\pagenumbering{arabic}\setcounter{page}{1}
\chapter{PENDAHULUAN}

\section{Latar Belakang}
Latar belakang tidak boleh terlalu panjang. Perhatikan pemenggalan kata, cetak miring, dan EYD. Pemenggalan kata di template ini enggak bisa otomatis sesuai ya, jadi kalau ada yang kurang sesuai, bisa disesuaikan secara manual. Ada 2 opsi, opsi satu ubah kata/kalimat, dan opsi dua dipenggal mandiri pakai "-".

\section{Pembatasan Masalah}
Isinya ya batasan masalah kamu, biar skripsi mu enggak merambat ke mana-mana.

\section{Tujuan dan Manfaat Penelitian}
Tujuan dan manfaat dari penulisan skripsi ini adalah sebagai berikut:
\begin{enumerate}
    \item Tujuan 1
    \item Tujuan 2
    \item Tujuan 3
\end{enumerate}

\section{Tinjauan Pustaka}
Berikan pengantar erlebih dahulu. Berikut merupakan beberapa literatur yang digunakan sebagai acuan dan referensi dalam penulisan skripsi ini. 


\section{Metodologi Penelitian}
Metodologi penelitian yang digunakan pada penulisan skripsi ini adalah studi literatur dan studi kasus. Studi literatur dilakukan dengan mengumpulkan informasi dari buku, jurnal, artikel, skripsi, \textit{paper}, dan sumber-sumber lainnya dari internet atau media cetak yang berkaitan dengan topik skripsi penulis. Informasi yang diperoleh digunakan sebagai acuan dan referensi dalam \ldots. Studi kasus dilakukan untuk \ldots. Algoritma yang digunakan pada penelitian ini adalah \ldots. Data yang digunakan sebagai studi kasus merupakan data sekunder berupa \ldots. Proses analisis dan komputasi dilakukan menggunakan software \ldots. 

\section{Sistematika Penulisan}
Pada penyusunan skripsi ini, penulis mengacu pada sistematika penulisan sebagai berikut:\\
\newline \textbf{BAB I PENDAHULUAN}
\par Bab ini berisi tentang latar belakang, pembatasan masalah, tujuan dan manfaat penelitian, tinjauan pustaka, metodologi penelitian, dan sistematika penulisan.\\
\newline \textbf{BAB II LANDASAN TEORI}
\par Bab ini berisi tentang dasar-dasar teori yang digunakan untuk membahas permasalahan yang ada dalam penelitian ini. Bab ini mencakup dasar teori tentang \ldots.\\
\newline \textbf{BAB III XXXXX}
\par Bab ini menjelaskan mengenai \ldots.\\
\newline \textbf{BAB IV STUDI KASUS}
\par Bab ini menjelaskan mengenai \ldots.\\
\newline \textbf{BAB V PENUTUP}
\par Bab ini berisi kesimpulan yang didapat berdasarkan hasil analisis dan saran yang dapat digunakan untuk penelitian selanjutnya.
